% Generated draft: requires mathematical and source review.
\documentclass[11pt]{article}
\usepackage[margin=25mm]{geometry}
\usepackage{fontspec}
\setmainfont{Latin Modern Roman}
\usepackage{amsmath,amssymb,mathtools,unicode-math}
\setmathfont{Latin Modern Math}
\usepackage{enumitem,needspace,xurl,hyperref}
\setlength{\parindent}{0pt}
\setlength{\parskip}{5pt}
\setlength{\emergencystretch}{3em}
\newcommand{\N}{\mathbb{N}}
\newcommand{\Z}{\mathbb{Z}}
\newcommand{\Q}{\mathbb{Q}}
\newcommand{\R}{\mathbb{R}}
\newcommand{\C}{\mathbb{C}}
\newcommand{\F}{\mathbb{F}}
\newcommand{\A}{\mathbb{A}}
\newcommand{\PP}{\mathbb{P}}
\newcommand{\E}{\mathbb{E}}
\newcommand{\eps}{\varepsilon}
\newcommand{\dd}{\,\mathrm{d}}
\begin{document}
\section*{UnsolvedMath ID: 700503}
\textbf{Source problem code:} catalogue entry
\subsection{Definitions and mathematical statement}
Let \(m\) denote the rank parameter in the source, and let \(F_m\) be the free group on \(m\) generators. For a positive integer \(n\), define the free Burnside group by
\[
B(m,n)=F_m/\langle\!\langle w^n : w\in F_m\rangle\!\rangle,
\]
where the denominator is the normal subgroup generated by all \(n\)-th powers. Thus \(B(m,n)\) is the free group in the class of groups satisfying \(x^n=1\) for every element \(x\). Write \(\operatorname{Aut}(G)\) for the group of bijective homomorphisms from \(G\) to itself, with composition as its operation.

\textbf{(a) Main question.} For sufficiently large odd positive integers \(n\), describe
\[
\operatorname{Aut}(B(m,n)).
\]
The source does not specify the permitted range of \(m\), a lower bound on \(n\), whether such a bound must be uniform in \(m\), or what form of description is required. These qualifications remain unresolved; the request is a structural description problem rather than a specified proposition to prove.

\textbf{(b) Archived question.} For sufficiently large odd positive integers \(n\), is the following assertion true?
\[
\forall H\leq B(m,n),\qquad
H\text{ noncyclic}\ \Longrightarrow\
\exists K\leq H\text{ such that }K\cong B(2,n).
\]
Here a group is cyclic if it is generated by one element. The rank range and the meaning of sufficiently large are again unspecified in the original question. The supplied archive reports an affirmative answer to part (b) for odd \(n\geq 1003\), attributed to V. S. Atabekyan. This is a supplied report of a solved part, not an independently verified theorem statement or an answer to part (a).
\subsection{Short English statement}
For sufficiently large odd exponents, describe the automorphisms of the free Burnside group \(B(m,n)\). A separate archived part asks whether every noncyclic subgroup contains a copy of \(B(2,n)\); the supplied record reports that part as solved affirmatively for odd \(n\geq 1003\). The rank range, the threshold for the main question, and the intended form of description require clarification.
\subsection{Sources}
\begin{itemize}
\item Supplied local UnsolvedMath registry record, permanent ID \(700503\), exact source code \texttt{catalogue entry}, titled \emph{Kourovka Notebook Problem 8.53}; upstream revision \texttt{d371ee917ac9647fbb06cd6fc529deec6f3dda82}.
\item A. Yu. Olshanskii, proposer of Problem 8.53 in \emph{The Kourovka Notebook (Unsolved Problems in Group Theory)}, 8th issue (1982). The supplied record identifies an update in the 21st issue (2026), \texttt{arXiv:1401.0300v48}. The collection description names E. I. Khukhro and V. D. Mazurov as editors.
\item The supplied Archive of Solved Problems entry attributes the affirmative result for part (b), for odd \(n\geq 1003\), to V. S. Atabekyan, \emph{Izv. Math.}, volume 73, no. 5 (2009), pp. 861--892. No article title or link is supplied.
\item The supplied registry status is \texttt{partially\_solved}. Its status review is dated \texttt{2026-10-16}, labels the evidence as published, and states that solved parts appear in the archive while remaining parts appear in the main section. This is supplied dated curation, not an independently verified literature assessment.
\end{itemize}
\end{document}
